%%% -*- coding: utf-8 -*-

\section{C\#}

\mode<beamer>{%
\showtitlepage{C\#}{C\#}
\lecture{C\#}
}

\subsection{Java vs.\ C\#}

\begin{frame}{Java vs.\ C\#: Big Picture}
  \begin{itemize}
  \item Both are class-based, object-oriented, statically typed languages
    with a similar-looking syntax (curly braces, semicolons, single
    inheritance of classes, interfaces).
  \item Both compile to an intermediate bytecode that runs on a virtual
    machine with a JIT compiler and garbage collection.
    \begin{itemize}
    \item Java: source $\to$ \texttt{.class} bytecode $\to$ JVM
    \item C\#: source $\to$ intermediate language (\texttt{.dll}/\texttt{.exe})
      $\to$ CLR (in Unity: Mono or IL2CPP)
    \end{itemize}
  \item Despite the similarity, several design decisions differ --
    this section highlights the ones you will notice first.
  \end{itemize}
\end{frame}

\begin{frame}{Naming Conventions}
  \begin{itemize}
  \item Java: \texttt{camelCase} for methods and fields,
    \texttt{PascalCase} for classes.
  \item C\#: \texttt{PascalCase} for classes, methods, \emph{and}
    properties; \texttt{camelCase} for local variables and parameters.
  \item Consequence: a method that reads a value is
    \texttt{getScore()} in Java, but simply \texttt{Score} (a property,
    see below) or \texttt{GetScore()} in C\#.
  \end{itemize}
\end{frame}

\begin{frame}[fragile]{Properties Instead of Getters/Setters}
  \begin{columns}[T]
    \column{0.5\textwidth}
      \textbf{Java}
      \begin{lstlisting}[language=Java]
private int score;

public int getScore() {
  return score;
}
public void setScore(int s) {
  score = s;
}
      \end{lstlisting}
    \column{0.5\textwidth}
      \textbf{C\#}
      \begin{lstlisting}
private int score;

public int Score {
  get { return score; }
  set { score = value; }
}

// or, auto-implemented:
public int Score { get; set; }
      \end{lstlisting}
  \end{columns}
  \begin{itemize}
  \item Properties look like fields to the caller (\texttt{obj.Score = 3;})
    but behave like methods.
  \end{itemize}
\end{frame}

\begin{frame}{Value Types vs.\ Reference Types}
  \begin{itemize}
  \item Java: primitives (\texttt{int}, \texttt{double}, \texttt{boolean},
    \ldots) are value types; \emph{everything else} is a reference type
    and lives on the heap.
  \item C\#: you can define your own value types with \texttt{struct}
    (copied by value, usually on the stack) in addition to
    \texttt{class} (reference type, on the heap).
  \item Relevant for Unity: \texttt{Vector3}, \texttt{Quaternion},
    \texttt{Color} are \texttt{struct}s -- assigning or passing them
    copies the whole value. A common bug when coming from Java is
    expecting mutations to a struct field of another object to
    "stick".
  \end{itemize}
\end{frame}

\begin{frame}{Nullability}
  \begin{itemize}
  \item Java: any reference may be \texttt{null}; primitives cannot be
    \texttt{null}.
  \item C\#: reference types may be \texttt{null} just like in Java;
    value types are \emph{not} nullable unless explicitly wrapped,
    e.g.\ \texttt{int?\ x = null;} (short for \texttt{Nullable<int>}).
  \item Convenience operators without a Java equivalent:
    \begin{itemize}
    \item \texttt{a?.Foo()} -- null-conditional call, evaluates to
      \texttt{null} instead of throwing if \texttt{a} is \texttt{null}.
    \item \texttt{a ?? b} -- null-coalescing, yields \texttt{b} if
      \texttt{a} is \texttt{null}.
    \end{itemize}
  \end{itemize}
\end{frame}

\begin{frame}{Namespaces and Packages}
  \begin{itemize}
  \item Java: package name must match the directory structure; one
    public top-level class per file; \texttt{import} for other packages.
  \item C\#: \texttt{namespace} is independent of the folder layout;
    a file may contain several classes and even several namespaces;
    \texttt{using} instead of \texttt{import}.
  \end{itemize}
\end{frame}

\begin{frame}{Access Modifiers}
  \begin{itemize}
  \item Java: \texttt{public}, \texttt{protected}, \texttt{private},
    and \emph{package-private} (the default -- no keyword).
  \item C\#: \texttt{public}, \texttt{protected}, \texttt{private},
    \texttt{internal} (visible within the whole assembly, not just a
    folder/package), plus combinations
    \texttt{protected internal} and \texttt{private protected}.
  \item C\# has no direct equivalent of Java's package-private; the
    default access (with no modifier) is \texttt{private} for members
    and \texttt{internal} for top-level types.
  \end{itemize}
\end{frame}

\begin{frame}{Delegates, Lambdas, and Events}
  \begin{itemize}
  \item Java: callbacks are modeled with functional interfaces
    (\texttt{Runnable}, \texttt{Comparator<T>}, \texttt{@FunctionalInterface})
    plus lambda expressions.
  \item C\#: \texttt{delegate} is a built-in, type-safe function-pointer
    type; \texttt{Action}/\texttt{Func<...>} are predefined generic
    delegates so you rarely declare your own.
  \item The \texttt{event} keyword builds the observer pattern into the
    language (subscribe with \texttt{+=}, unsubscribe with \texttt{-=}) --
    Java has no direct counterpart.
  \end{itemize}
\end{frame}

\begin{frame}[fragile]{Operator Overloading}
  \begin{itemize}
  \item Java: not allowed (the \texttt{+} operator on \texttt{String}
    is a special case built into the language).
  \item C\#: user-defined types can overload operators -- used
    extensively by Unity's math types.
  \end{itemize}
  \begin{lstlisting}
public static Vector3 operator +(Vector3 a, Vector3 b)
  => new Vector3(a.x + b.x, a.y + b.y, a.z + b.z);
  \end{lstlisting}
\end{frame}

\begin{frame}{Exceptions}
  \begin{itemize}
  \item Java distinguishes \emph{checked} exceptions (must be declared
    with \texttt{throws} or caught) from \emph{unchecked} ones
    (\texttt{RuntimeException} and subclasses).
  \item C\# has only unchecked exceptions -- no \texttt{throws}
    declaration exists, and the compiler never forces you to catch
    anything.
  \end{itemize}
\end{frame}

\begin{frame}[fragile]{Collections: Streams vs.\ LINQ}
  \begin{columns}[T]
    \column{0.5\textwidth}
      \textbf{Java Streams}
      \begin{lstlisting}[language=Java]
names.stream()
  .filter(n -> n.length() > 3)
  .map(String::toUpperCase)
  .collect(Collectors.toList());
      \end{lstlisting}
    \column{0.5\textwidth}
      \textbf{C\# LINQ}
      \begin{lstlisting}
names
  .Where(n => n.Length > 3)
  .Select(n => n.ToUpper())
  .ToList();
      \end{lstlisting}
  \end{columns}
\end{frame}

\begin{frame}[fragile]{Small Syntax Differences}
  \begin{itemize}
  \item String interpolation: \texttt{\$"Hello, \{name\}!"} in C\#
    vs.\ \texttt{"Hello, " + name + "!"} or \texttt{String.format(...)}
    in Java.
  \item Both languages have \texttt{var} for local type inference; C\#
    additionally allows target-typed \texttt{new()}, e.g.\
    \texttt{List<int> xs = new();}
  \item C\# supports \emph{extension methods} -- adding a method to an
    existing type without subclassing or modifying it. Java has no
    equivalent (default interface methods come closest but require an
    interface).
  \item Generics: C\# generics are reified (the type is known at run
    time); Java generics use type erasure.
  \end{itemize}
\end{frame}

\begin{frame}{Takeaways for Unity}
  \begin{itemize}
  \item \texttt{MonoBehaviour} callbacks (\texttt{Start}, \texttt{Update},
    \ldots) are ordinary C\# methods matched by name and signature --
    no \texttt{@Override} annotation, so a typo silently creates an
    unused method instead of a compiler error.
  \item Math and value types (\texttt{Vector3}, \texttt{Quaternion},
    \texttt{Color}, \texttt{struct}s in general) are copied by value --
    watch out for the Java habit of expecting shared, mutable state.
  \item Coroutines (\texttt{IEnumerator} + \texttt{yield return}) provide
    cooperative multitasking with no direct Java counterpart.
  \end{itemize}
\end{frame}

%%% Local Variables:
%%% mode: latex
%%% TeX-master: "main-slides"
%%% mode: reftex
%%% mode: flyspell
%%% End:
